\chapter{Alimentazioni e accensione}
\label{ch:power}

\begin{semplice}
La scheda prende i \SI{5}{\volt} dal dock del satellite oppure dalla USB-C; se ci sono tutti e due vince il dock. Dai \SI{5}{\volt} un convertitore ricava i \SI{3.3}{\volt} (il \enquote{3V3}) che alimentano quasi tutto. Il convertitore però si accende solo se il dock gli dà il permesso con il segnale EN\_DOCK: con la sola USB-C la scheda resta spenta.
\end{semplice}

\section{Rail di alimentazione}
Un \enquote{rail} è una linea di alimentazione. Nella tabella: \textbf{Rail} = nome della net; \textbf{Generato da} = il componente o il connettore che la fornisce; \textbf{Valore} = tensione nominale ed eventuale intervallo garantito; \textbf{Carichi principali} = cosa alimenta.

\begin{xltabular}{\linewidth}{@{}L{24mm}L{30mm}L{30mm}Y@{}}
\tabhead\hd{Rail}&\hd{Generato da}&\hd{Valore}&\hd{Carichi principali}\\
\endhead
\net{5V} & dock J708 pin 1 e 3 & \SI{5}{\volt} nominali & ingresso IN1 di IC4 (priorità)\\
\net{/POWER/VBUS} & USB-C J15 & \SI{5}{\volt}. Le resistenze da \SI{5.1}{\kilo\ohm} sui pin CC1/CC2 dicono al PC che la scheda è un dispositivo USB 2.0 & ingresso IN2 di IC4\\
\net{5V0} & IC4 TPS2121 (OUT) & uguale all'ingresso scelto & IC3 (buck), VDD di IC202 (switch RF)\\
\net{3V3} & IC3 TPSM863257 & \SI{3.30}{\volt} nominali, \SIrange{3.25}{3.405}{\volt} \fonte{REV-1104 §2.1} & tutta la logica, Si4464, TCXO, PA, LNA, sensori\\
\net{RTC\_VBAT} & 3V3 via D401 (BAT54WS) e R412 \SI{1}{\kilo\ohm} & $\approx$\SI{3.15}{\volt} \fonte{SCH nota MCU} & VBAT di U404, 10 supercondensatori CPH3225A (\SI{110}{\milli\farad})\\
\bottomrule
\end{xltabular}

Il pin VBAT dell'MCU (pin 1) e VDDUSB (pin 48) sono collegati direttamente al 3V3 \fonte{SCH}. Il dominio di backup dell'STM32 non ha una batteria propria: l'RTC interno dell'MCU perde l'ora quando il 3V3 si spegne. Il riferimento di tempo è il DS3231M (capitolo~\ref{ch:supervisione}).

I pin \net{+12V}, \net{+3V3\_CONN}, \net{RS485.A1/B1} del dock non sono usati su questa scheda \fonte{LOG R-12; SCH nota IF}.

\section{Power mux TPS2121: priorità al dock}
\begin{itemize}
\item \textbf{Priorità IN1 (dock).} PR1 = $5 \times 4{,}7/(10+4{,}7) = \SI{1.60}{\volt}$, sopra $V_\text{REF} = \SI{1.06}{\volt}$: con il dock presente vince sempre il dock \fonte{LOG D1; REV-1104 §2.1}. Il dock perde la priorità quando scende sotto $1{,}06 \times 14{,}7/4{,}7 = \SI{3.32}{\volt}$.
\item \textbf{Sovratensione IN1.} Soglia OV1 = $1{,}06 \times (22+4{,}7)/4{,}7 = \SI{6.02}{\volt}$ nominali \fonte{LOG rev.\ 1.44, P-08}. OV2 (USB) è disattivata (a GND).
\item \textbf{Limite di corrente.} ILIM con R44 = \SI{33}{\kilo\ohm}: \SI{3.21}{\ampere} \fonte{REV-1104 §2.1}.
\item \textbf{Stato.} Il pin ST è a GND: \textbf{il firmware non sa quale ingresso è attivo}.
\end{itemize}

\section{Buck TPSM863257 e segnale EN\_DOCK}
L'ingresso EN di IC3 è pilotato da \net{/IF/EN\_DOCK} (J708 pin 15) attraverso R46 \SI{10}{\kilo\ohm}. EN ha un pull-down interno da \SI{2}{\mega\ohm}: \textbf{senza EN\_DOCK alto il 3V3 resta spento} \fonte{LOG Parte B; REV-1104 §2.1}. È una scelta di progetto confermata \fonte{LOG D1}.

Il soft-start interno dura \SI{1.6}{\milli\second} tipici \fonte{DS-TPSM863257 p.\,7}. L'uscita PG non è collegata.

\begin{attenzione}[alimentazione al banco]
Con la sola USB-C il 5V0 c'è, ma il 3V3 no: l'MCU, il CH340E e la console sono spenti. Al banco bisogna fornire anche EN\_DOCK (EN max raccomandato \SI{5.5}{\volt}) \fonte{REV-1104 §2.1}. In trasmissione il 5V0 assorbe circa \SI{1}{\ampere}: una porta USB 2.0 di default dà \SI{500}{\milli\ampere}. Usare il dock o un alimentatore da almeno \SI{1.5}{\ampere} \fonte{REV-1104 P22}.
\end{attenzione}

\section{Che cosa vede il firmware}
L'hardware non porta all'MCU nessun segnale di stato dell'alimentazione (né ST del TPS2121, né PG del buck, né una misura del 5\,V). Il firmware ha a disposizione:
\begin{itemize}
\item \textbf{VDDA tramite VREFINT} (ADC1 canale 0):
\[ V_\text{DDA} = \SI{3.0}{\volt}\times\frac{\text{VREFINT\_CAL}}{\text{VREFINT\_DATA}} \]
VREFINT\_CAL si legge all'indirizzo \code{0x0BFA05AA}; è misurato in fabbrica a \SI{3.0}{\volt} e \SI{30}{\celsius} \fonte{DS-STM32L562 p.\,65}. Il 3V3 è anche VDDA, quindi questa è la misura del rail 3V3.
\item \textbf{Causa del reset} in \code{RCC\_CSR} (BOR, pin NRST, IWDG\ldots) \fonte{LOG F-13}.
\item \textbf{Pin RST del DS3231M} su PB9: va basso quando il VCC dell'RTC scende sotto la soglia di power-fail $V_\text{PF}$ \fonte{DS-DS3231M p.\,7}. Il valore numerico di $V_\text{PF}$ non è stato estratto dal datasheet: \DV.
\item \textbf{Temperatura del buck} (PT1000 su PC1) e del PA (PT1000 su PB1).
\item \textbf{Risposta del Si4464} dopo ogni TX: se ha perso la configurazione c'è stato un calo del 3V3 (nota F-10, paragrafo~\ref{sec:rf-supervisore}).
\end{itemize}

\section{Sequenza di accensione}
\label{sec:powerup}
La figura~\ref{fig:powerup} riassume cosa succede dal momento in cui il dock alza EN\_DOCK. I tempi sono tipici; i valori minimi e massimi sono nella tabella~\ref{tab:powerup}.

\begin{figure}[!htb]
\centering
\begin{tikzpicture}[x=0.72mm,y=1mm,font=\sffamily\scriptsize,
   sig/.style={thick},dim/.style={{Stealth[length=1.6mm]}-{Stealth[length=1.6mm]},sbgray}]
% righe: base (livello basso) a y = -13*i, livello alto a +5
\foreach \i/\name/\sub in {0/EN\_DOCK/dal dock,1/3V3/uscita del buck,2/NRST/reset MCU,3/WDI/impulsi del firmware,4/SDN/alto = radio spenta (modo A),5/RST DS3231M/avviso dell'RTC,6/WDO/reset dal watchdog}{
   \node[anchor=east,align=right] at (-2,-13*\i+2.5) {\textbf{\name}\\[-1pt]{\color{sbgray}\sub}};
   \draw[sbgray!30] (0,-13*\i) -- (150,-13*\i);}
% linee verticali degli eventi
\foreach \x/\t in {8/{A},14/{B},30/{C},72/{D}}{
   \draw[sbgray!70,densely dashed] (\x,8) -- (\x,-82);
   \node[circle,draw=sbgray,inner sep=0.8pt,fill=white] at (\x,10) {\t};}
% EN_DOCK
\draw[sig,sbgreen] (0,0) -- (8,0) -- (8,5) -- (150,5);
% 3V3
\draw[sig,sbred] (0,-13) -- (10,-13) -- (14,-8) -- (150,-8);
\draw[dim] (8,-15.5) -- (14,-15.5); \node[anchor=west,fill=white,inner sep=1pt] at (15,-15.5) {soft-start \SI{1.6}{\milli\second}};
% NRST
\draw[sig,sbblue] (0,-26) -- (14,-26) -- (14,-21) -- (150,-21);
% WDI
\draw[sig,sbblue] (0,-39) -- (30,-39) -- (30,-34) -- (60,-34) -- (60,-39) -- (90,-39) -- (90,-34) -- (120,-34) -- (120,-39) -- (150,-39);
\draw[dim] (14,-41.5) -- (30,-41.5); \node[anchor=west,fill=white,inner sep=1pt] at (31,-41.5) {primo fronte entro \SI{1.12}{\second}};
% SDN
\draw[sig,sborange] (0,-52) -- (14,-52) -- (14,-47) -- (72,-47) -- (72,-52) -- (150,-52);
\draw[dim] (14,-54.5) -- (72,-54.5); \node[anchor=west,fill=white,inner sep=1pt] at (73,-54.5) {$t_d$ = \SIrange{120}{300}{\milli\second}};
% RST DS3231M
\draw[sig,sbgreen] (0,-65) -- (78,-65) -- (78,-60) -- (150,-60);
\draw[dim] (14,-67.5) -- (78,-67.5); \node[anchor=west,fill=white,inner sep=1pt] at (79,-67.5) {$t_\text{REC}$ fino a \SI{300}{\milli\second}};
% WDO
\draw[sig,sbgreen] (0,-78) -- (14,-78) -- (14,-73) -- (150,-73);
\node[anchor=north,inner sep=1pt] at (118,-73.4) {resta alto: il firmware manda gli impulsi};
\end{tikzpicture}
\caption{Sequenza di accensione (schema non in scala). A: il dock alza EN\_DOCK. B: il 3V3 è pronto e l'MCU parte. C: primo impulso del firmware al watchdog. D: il supervisore libera la radio. Dopo C il firmware deve cambiare il livello di WDI almeno una volta al secondo. Linea bassa = livello logico 0, linea alta = livello logico 1.}
\label{fig:powerup}
\end{figure}

\begin{xltabular}{\linewidth}{@{}L{8mm}L{48mm}L{36mm}Y@{}}
\caption{Eventi all'accensione e tempi con fonte.}\label{tab:powerup}\\
\tabhead\hd{\#}&\hd{Evento}&\hd{Tempo}&\hd{Fonte e conseguenza per il firmware}\\
\endfirsthead
\tabhead\hd{\#}&\hd{Evento}&\hd{Tempo}&\hd{Fonte e conseguenza per il firmware}\\
\endhead
1 & Soft-start del buck, 3V3 sale & \SI{1.6}{\milli\second} tipici & \fonte{DS-TPSM863257 p.\,7}\\
2 & POR/BOR dell'MCU, boot da flash (BOOT0 = L tramite R3 \SI{10}{\kilo\ohm} a GND) & \textendash & \fonte{SCH}\\
3 & IC205 tiene SDN del Si4464 alto dopo che il 3V3 supera la soglia & $t_d$ = \SI{120}{}/\SI{200}{}/\SI{300}{\milli\second} (min/tip/max) & \fonte{DS-TPS3825 p.\,8}. Non usare il Si4464 prima di $t_d$ max: vedi F-05.\\
4 & Soglia del supervisore IC205 & \SI{2.88}{}/\SI{2.93}{}/\SI{3.00}{\volt} & \fonte{LOG §46; DS-TPS3825 p.\,6}\\
5 & Il Si4464 esce dallo shutdown: POR e calibrazioni interne, poi serve POWER\_UP & Shutdown$\to$TX o RX: \SI{15}{\milli\second} (indicativo) & \fonte{DS-Si4464 p.\,20--21}\\
6 & RST del DS3231M basso dopo il power-up & $t_\text{REC}$ = \SI{250}{\milli\second} tip., \SI{300}{\milli\second} max & \fonte{DS-DS3231M p.\,3}. Non influisce su I2C e RTC \fonte{DS-DS3231M p.\,4}.\\
7 & Prima scrittura possibile sulla FRAM & $t_\text{PU}$ $\geq$ \SI{250}{\micro\second} & \fonte{DS-FM25V10 p.\,19}\\
8 & Primo impulso su WDI (PC6) & entro \SI{1.12}{\second} ($t_{WD}$ min) & \fonte{LOG F-12; DS-STWD100}\\
9 & Il 5V0 arriva al VDD dello switch prima dei comandi a 3V3 & sequenza rispettata dall'hardware & \fonte{DS-HMC349A p.\,8; REV-1104 §2.2}\\
\bottomrule
\end{xltabular}

Nella tabella~\ref{tab:powerup} gli eventi sono in ordine di tempo. I valori \enquote{min/tip/max} sono minimo, tipico e massimo garantiti dal datasheet: il firmware deve sempre usare il valore peggiore (il massimo per le attese, il minimo per le scadenze). POR = reset automatico all'accensione; BOR = reset per tensione troppo bassa.

\begin{nota}[stato sicuro durante il boot]
Durante il reset e il boot dell'MCU i GPIO sono in alta impedenza. Il pull-up R702 tiene alto il gate di Q701 e la net di PC2: switch su RX, APC a zero (PA spento), LNA acceso. La radio è sotto il controllo di IC205 (modo A) \fonte{LOG rev.\ 1.41; REV-1104 §3}. Lo stato di default è quindi sicuro anche senza firmware.
\end{nota}
